% UFC Fight Winner Prediction — TMA4268 Compulsory Exercise 2

\documentclass{article}
\usepackage[T1]{fontenc}
\usepackage[utf8]{inputenc}
\usepackage{graphicx}
\usepackage{booktabs}
\usepackage{amsmath}
\usepackage[protrusion=true,expansion=false]{microtype}
\usepackage{url}
\usepackage{hyperref}
\usepackage[margin=1in]{geometry}
\usepackage{caption}
\usepackage{subcaption}
\usepackage{float}
\emergencystretch=3em

\title{Predicting UFC Fight Outcomes Using Pre-Fight Fighter Statistics\\
\large Compulsory Exercise 2 --- TMA4268 Statistical Learning V2026}
\author{Ricardo Sonda Guiraudeli; Vegard Ragnhildstveit; Anders Bekkevard}
\date{April 2026}

\begin{document}

\maketitle

%% -------------------------------------------------------------------------
\section*{Abstract}
%% -------------------------------------------------------------------------

We examine the feasibility of predicting the winner of a UFC martial arts
fight from pre-fight statistics only.
The target is a binary outcome: whether the fighter assigned to the Red
corner wins (\texttt{WinnerRed} = 1) or loses (\texttt{WinnerRed} = 0).
The dataset is the \emph{Ultimate UFC Dataset} sourced from Kaggle
(\texttt{mdabbert/ultimate-ufc-dataset}), containing 6{,}528 bouts between
March 2010 and December 2024 with per-fighter cumulative statistics recorded
at fight time.

We make use of a \emph{difference features} feature set - each numeric statistic is
replaced by its Red-minus-Blue signed difference - to remove redundancy 
in the dataset and reduce dimensionality from 118 to 26 predictors.
Three model families are compared: penalised logistic regression (Elastic
Net, tuned with nested cross-validation), a Random Forest ensemble, and
gradient-boosted trees (XGBoost).
All models are evaluated on a chronological 80/20 holdout split to prevent
time leakage, and benchmarked against the implied win probabilities
derived from pre-fight betting odds.

The bookmaker (betting-implied) baseline achieves an accuracy of 69.6\% and a ROC AUC of 0.743.
The model with best performance, the Random Forest, reaches an accuracy of 63.6\%
and a ROC AUC of 0.693.
The gap to the bookmaker confirms that publicly available fighter statistics
contain real predictive signal but fall well short of the full information
encoded in market odds, as expected.
Among observable predictors, the average takedown-landed difference, reach
advantage, and win-streak difference emerge as the strongest predictors of
fight outcome, while a negative age coefficient reveals that an age gap is a 
significant disadvantage for older fighters.

%% -------------------------------------------------------------------------
\section{Introduction}
%% -------------------------------------------------------------------------

Mixed martial arts (MMA) is one of the fastest-growing combat sports
globally, and the UFC is its premier organisation.
Fights can be thought of is a highly complex event influenced by physical attributes,
technical skills, recent form, and matchup dynamics.
Predicting fight outcomes is a classic classification problem: given
observable pre-fight information about both fighters, can we assign a
probability that the Red-corner fighter wins?

\paragraph{Dataset.}
We use the \emph{Ultimate UFC Dataset} published by user \texttt{mdabbert}
on Kaggle (\texttt{mdabbert/ultimate-ufc-dataset}).
The dataset records 6{,}528 bouts with 118 raw columns, including
per-fighter cumulative career statistics (wins, losses, streaks, striking
and grappling averages, physical measurements), contextual fight information
(weight class, number of rounds, gender), and bookmaker moneyline odds.
After filtering to bouts with a decisive Red or Blue winner (no ties) and dropping
the pre-computed difference columns provided by the source, we retain 6{,}275
bouts for analysis.

\paragraph{Task and scope.}
Our goal is to build a classifier that, given only pre-fight information,
assigns a calibrated probability of a Red-corner victory.
We benchmark every model against bookmaker-implied probabilities --- a
pretty-much-near-market-efficient baseline that absorbs information beyond public
statistics.
This project is meant as an exercise in the statistical methods employed in the
statistical learning course and we do not intend to create a model that beats 
the moneyline odds. 

\paragraph{Practical constraint.}
Because leakage from betting odds onto odds-derived features would trivially dominate any model, 
odds-derived columns are excluded from all feature sets; the odds appear only as a
benchmark in evaluation.

\paragraph{How moneyline odds work.}
UFC betting odds are usually quoted as \emph{American moneyline} odds.
A negative line (e.g., $-150$) indicates the favourite: you must stake 150
units to win 100 units of profit.
A positive line (e.g., $+130$) indicates the underdog: a 100-unit stake wins
130 units of profit.
These lines can be converted to implied win probabilities using
$p = \frac{|o|}{|o| + 100}$ for negative odds $o < 0$, and
$p = \frac{100}{o + 100}$ for positive odds $o > 0$.
In practice, bookmakers include a margin (``vig''), so the two implied
probabilities in a matchup typically sum to slightly more than 1.
In this project, we generally work with implied probabilities instead of the 
moneyline odds.

%% -------------------------------------------------------------------------
\section{Descriptive Data Analysis}
%% -------------------------------------------------------------------------

\subsection{Dataset overview}

After cleaning, the dataset contains 6{,}275 bouts spanning March 2010 to
December 2024.
The outcome is moderately imbalanced: 58.1\% of bouts result in a Red-corner
victory, as shown in Figure~\ref{fig:class_balance}.
This imbalance reflects the UFC's convention of assigning the Red corner to
champions and higher-ranked fighters, who naturally win more often.
Table~\ref{tab:summary_stats} summarises the key numeric predictors used in
modelling.

\begin{figure}[H]
\centering
\includegraphics[width=0.9\textwidth]{../figures/eda/class_balance_timeline.png}
\caption{Left: Red vs Blue win counts in the cleaned dataset (6{,}275 bouts).
Red-corner fighters win 58.1\% of bouts.
Right: number of bouts per year; data coverage is near-complete from 2013
onwards.}
\label{fig:class_balance}
\end{figure}

\begin{table}[H]
\centering
\caption{Summary statistics for selected difference features
(\emph{Red minus Blue}) on the full cleaned dataset ($n = 6{,}275$).
Missing-value counts differ per feature because career statistics were
not tracked for early bouts.}
\label{tab:summary_stats}
\begin{tabular}{lrrrr}
\toprule
Feature & Mean & Std & Min & Max \\
\midrule
\texttt{AgeDiff}               &  0.56 & 5.18 & $-16.0$ & 17.0 \\
\texttt{ReachCmsDiff}          &  0.23 & 8.62 & $-30.5$ & 188.0 \\
\texttt{CurrentWinStreakDiff}  &  0.14 & 1.87 & $-10.0$ & 18.0 \\
\texttt{CurrentLoseStreakDiff} &  0.12 & 1.01 & $-5.0$  & 6.0 \\
\texttt{AvgSigStrLandedDiff}   &  0.32 & 17.4 & $-128.2$ & 104.7 \\
\texttt{AvgSigStrPctDiff}      &  0.01 & 0.14 & $-0.57$ & 0.82 \\
\texttt{AvgTDLandedDiff}       &  0.09 & 1.79 & $-10.9$ & 11.0 \\
\texttt{WinsDiff}              &  1.49 & 4.17 & $-23.0$ & 28.0 \\
\bottomrule
\end{tabular}
\end{table}

The near-zero means for the difference features are expected, since the UFC
broadly matches fighters of comparable ability, so the average edge is
close to zero.
The positive mean for \texttt{WinsDiff} (1.49) and \texttt{AgeDiff} (0.56)
reflects the aforementioned policy for assigning more experienced fighters
to the red corner.

\subsection{Missingness}

Many statistics columns have substantial missing values for early
bouts during when tracking was incomplete. Also, a substantial number of features
related to weight-classes are empty both due to the sparse encoding and due unrecorded
data for unranked fighters. 
Figure~\ref{fig:missing_bar} shows the completeness rate per column.
The \emph{diff} feature set strategy (Section~\ref{sec:methods}) retains
all rows and delegates imputation to the model pipeline.

\begin{figure}[H]
\centering
\includegraphics[width=0.85\textwidth]{../figures/missing/missing_bar.png}
\caption{Fraction of non-missing values per column in the raw dataset.
Fighter striking and grappling averages have notably lower completeness
for older bouts.}
\label{fig:missing_bar}
\end{figure}

\subsection{Correlation with outcome}

Figure~\ref{fig:corr_target} shows the Pearson correlation of each numeric
raw column with the binary outcome \texttt{WinnerRed}.
Correlations are modest ($|r| < 0.12$), confirming that no single feature
is strongly predictive on its own.
Age, win streak, cumulative wins, and significant-strike statistics show the
highest marginal correlations.

\begin{figure}[H]
\centering
\includegraphics[width=0.75\textwidth]{../figures/correlation/corr_with_winner.png}
\caption{Pearson correlation of raw numeric columns with \texttt{WinnerRed}.
The strongest individual predictors are win-streak and total-wins
differences (Red corner), all with $|r| < 0.12$.}
\label{fig:corr_target}
\end{figure}

\subsection{Feature distributions by outcome}

Figure~\ref{fig:violin} shows violin plots of the six most important
difference features split by fight outcome.
The distributions largely overlap, which is consistent with the low
individual correlations.
The signal is subtle: in Red-win fights the win-streak and reach
distributions shift slightly upward, confirming a real but weak effect.

\begin{figure}[H]
\centering
\includegraphics[width=\textwidth]{../figures/eda/feature_violin_by_outcome.png}
\caption{Distribution of key difference features (Red minus Blue) split
by outcome. Dashed lines mark zero (even matchup). Distributions overlap
substantially, reflecting the difficulty of the prediction task.}
\label{fig:violin}
\end{figure}

\subsection{Feature correlation structure}

The fighter statistics are highly internally correlated (e.g.,
wins is correlated with win streak and total rounds fought).
This motivates using difference features to compress the Red/Blue
column pairs into a single signed quantity, reducing some of the multicolinearity
in the features.

%% -------------------------------------------------------------------------
\section{Methods}
\label{sec:methods}
%% -------------------------------------------------------------------------

\subsection{Feature engineering}

The raw dataset contains symmetrically paired Red/Blue columns for every
fighter statistic.
We replace each pair $(x_{\text{red}}, x_{\text{blue}})$ with the signed
difference $\Delta x = x_{\text{red}} - x_{\text{blue}}$.
This yields symmetric features: swapping Red and Blue
negates all difference features and flips the outcome, so the representation
is not biased on the corners.
The resulting \emph{diff} feature set contains 26 predictors (numeric
differences plus categorical context variables: gender, title-bout indicator,
and number of scheduled rounds).

Missing values are imputed inside using medians
for numeric features and mode for categorical features, so all 6{,}275 bouts
are preserved.

\subsection{Data split}

All models use a \textbf{chronological holdout split}: fights are sorted by
date, the earliest 80\% form the training set (5{,}020 bouts), and the most
recent 20\% form the test set (1{,}255 bouts; fights from approximately
late 2021 onwards).
This choice prevents \emph{time leakage}: since some numeric features are
cumulative fighter statistics, information about the results of past fights
are contained in the later entries. Using this information to predict
previous fights thus leads to leakage.

\subsection{Model 1: Logistic Regression (baseline)}

Logistic regression models the log-odds of Red winning as a linear function
of the (min-max scaled and imputed) difference features:
\[
  \log \frac{P(\text{Red wins})}{1 - P(\text{Red wins})} = \beta_0 + \boldsymbol{\beta}^\top \Delta\mathbf{x}.
\]
\textbf{Strengths}: interpretable coefficients, fast to fit, simple.
\textbf{Weaknesses}: assumes a linear decision boundary; cannot capture
non-linear interactions between features.
In our setting the feature space is noisy and low-signal, so a linear boundary
is a reasonable starting assumption.
No regularisation is applied for the baseline model. Coefficients are estimated
by likelihood.

\subsection{Model 2: Elastic Net Logistic Regression (penalised, tuned)}

An Elastic Net is just a Logistic Regression 
with hybrid $L_1$ and $L_2$ regularisation penalty to the log-likelihood:
\[
  \mathcal{L}(\boldsymbol{\beta}) - \lambda \left[
    \frac{1-\alpha}{2}\|\boldsymbol{\beta}\|_2^2
    + \alpha \|\boldsymbol{\beta}\|_1
  \right],
\]
where $C = 1/\lambda$ controls regularisation strength and $\alpha$
(called \texttt{l1\_ratio} in scikit-learn) interpolates between Ridge
($\alpha=0$) and Lasso ($\alpha=1$).
Elastic Net is preferable to pure Lasso when predictors are correlated:
Lasso tends to arbitrarily discard one of a correlated group, whereas
Ridge-like shrinkage stabilises the selection.

Our \textbf{hyperparameter tuning} uses a nested cross-validation structure:
\begin{itemize}
  \item \textbf{Outer split}: the chronological 80/20 holdout (shared with
  all models).
  \item \textbf{Inner CV}: 5-fold stratified cross-validation over the
  training set.
  \item \textbf{Grid}: $C \in \{0.01, 0.1, 1.0, 10.0\}$,
  $\alpha \in \{0.1, 0.5, 0.9, 1.0\}$ (16 combinations).
\end{itemize}
The best parameters found were $C = 0.01$ (strong regularisation),
$\alpha = 0.1$ (mostly Ridge with a light Lasso component).
21 out of 26 features survived the Lasso non-zero
coefficients.

\subsection{Model 3: Random Forest}

A Random Forest trains $B = 400$ decision trees on bootstrap samples and
averages their estimates.
Each tree considers $\sqrt{p}$ randomly selected features at each split.
\textbf{Strengths}: captures non-linear structure and feature interactions
without explicit specification; robust to outliers; provides feature
importances via mean decrease in impurity.
\textbf{Weaknesses}: less interpretable than linear models; can overfit with
deep trees on small datasets.
We set \texttt{min\_samples\_leaf = 10} to regularise tree depth and reduce
variance; this was chosen by our judgment given the dataset size
(5{,}020 training observations, 26 features).

\subsection{Model 4: Gradient Boosted Trees (XGBoost)}

XGBoost builds trees sequentially, with each tree correcting the residuals
of the current ensemble.
The learning rate (\texttt{learning\_rate} $= 0.05$) limits each tree's
contribution to prevent overfitting, and stochastic subsampling
(\texttt{subsample} $= 0.8$, \texttt{colsample\_bytree} $= 0.8$) introduces
variance reduction.
The minimum child weight (\texttt{min\_child\_weight} $= 5$) acts as
additional regularisation on leaf size.
\textbf{Strengths}: strong empirical performance; handles missing values
natively; flexible regularisation.
\textbf{Weaknesses}: more hyperparameters to configure; prone to overfitting
without careful tuning; less interpretable than logistic regression.

\subsection{Evaluation metrics}

All models are evaluated on the held-out chronological test set and compared
against the bookmaker-implied probabilities.
We report:
\begin{itemize}
  \item \textbf{Accuracy}: fraction of correctly predicted outcomes
  (threshold at 0.5).
  \item \textbf{Log loss} (cross-entropy): measures calibration and
  sharpness of probability estimates. Lower is better.
  Treated as the primary ranking metric because it rewards well-calibrated
  probabilities, not just the sign of the prediction.
  \item \textbf{Brier score}: mean squared error between predicted
  probabilities and binary outcomes. Lower is better.
  \item \textbf{ROC AUC}: area under the receiver operating characteristic
  curve. Measures the discrimination ability independent of the decision
  threshold.
\end{itemize}

%% -------------------------------------------------------------------------
\section{Results and Interpretation}
\label{sec:results}
%% -------------------------------------------------------------------------

\subsection{Model performance}

Table~\ref{tab:results} compares all models on the 1{,}255-bout test set.
The bookmaker benchmark is shown for reference. $\Delta$ values are
signed differences relative to the bookmaker (for negative $\Delta$ log
loss, Brier is better. For positive $\Delta$ accuracy, AUC is better).

\begin{table}[H]
\centering
\caption{Test-set evaluation results (chronological holdout, $n_{\text{test}} = 1{,}255$).}
\label{tab:results}
\begin{tabular}{lcccccc}
\toprule
Model & Acc & $\Delta$ Acc & Log Loss & $\Delta$ LL & Brier & ROC AUC \\
\midrule
Bookmaker (baseline)    & 0.696 & ---    & 0.593 & ---    & 0.204 & 0.743 \\
Naive (majority class)  & 0.571 & $-$0.124 & 0.693 & $+$0.100 & 0.250 & 0.500 \\
Logistic Regression     & 0.635 & $-$0.060 & 0.636 & $+$0.043 & 0.223 & 0.673 \\
Elastic Net (nested CV) & 0.640 & $-$0.056 & 0.635 & $+$0.041 & 0.222 & 0.682 \\
\textbf{Random Forest}  & \textbf{0.636} & $-$0.060 & \textbf{0.633} & $+$0.040 & \textbf{0.221} & \textbf{0.693} \\
XGBoost                 & 0.632 & $-$0.064 & 0.636 & $+$0.043 & 0.223 & 0.670 \\
\bottomrule
\end{tabular}
\end{table}

Random Forest achieves the best log loss (0.633), Brier score (0.221), and
ROC AUC (0.693).
However, the differences among Logistic Regression, Elastic Net, and Random
Forest are small (log loss range $< 0.004$), suggesting that the dominant
bottleneck is the limited content of the public statistics, not the
model family.
XGBoost performs slightly worse despite its higher complexity.

All models comfortably beat the naive majority-class baseline (AUC 0.5)
but remain well below the bookmaker (AUC 0.743).
One way to interpret the persistent $\approx$5.5 percentage-point accuracy gap
to the bookmaker is that it quantifies the information advantage of market 
prices over publicly available fight statistics.

\subsection{Confusion matrices}

Table~\ref{tab:cm} shows confusion matrix entries and precision/recall for
each model.
A common pattern is that all models predict Red wins the large majority
of the time, driven by the 58.1\% class imbalance.
Recall for the Red class (fraction of actual Red wins correctly called) is
above 0.83 for all linear/ensemble models, while recall for the Blue class
is much lower.
XGBoost achieves the most balanced predictions (TN = 233, highest among all
models), reflecting the effect of \texttt{min\_child\_weight} regularisation.

\begin{table}[H]
\centering
\caption{Confusion matrix entries and precision/recall on the test set
($n = 1{,}255$; 717 Red wins, 538 Blue wins).
Precision and recall are reported for the Red-wins class (positive class).}
\label{tab:cm}
\begin{tabular}{lrrrrrr}
\toprule
Model & TN & FP & FN & TP & Precision & Recall \\
\midrule
Logistic Regression     & 202 & 336 & 122 & 595 & 0.639 & 0.830 \\
Elastic Net (nested CV) & 198 & 340 & 112 & 605 & 0.640 & 0.844 \\
Random Forest           & 201 & 337 & 120 & 597 & 0.639 & 0.833 \\
XGBoost                 & 233 & 305 & 157 & 560 & 0.647 & 0.781 \\
\bottomrule
\end{tabular}
\end{table}

\subsection{ROC curves}

Figure~\ref{fig:roc} shows the ROC curve for the best-performing model
(Random Forest) alongside the bookmaker benchmark.
The bookmaker ROC curve lies clearly above the model curve across all
thresholds, confirming that market odds encode more discriminative
information than our feature set.

\begin{figure}[H]
\centering
\includegraphics[width=0.58\textwidth]{../figures/evaluation/random_forest_diff_roc.png}
\caption{ROC curves for the Random Forest model versus the bookmaker
baseline on the test set. The bookmaker (AUC = 0.743)
outperforms the best statistical model (AUC = 0.693).}
\label{fig:roc}
\end{figure}

\subsection{Coefficient interpretation (Elastic Net)}

Table~\ref{tab:coefs} shows the largest-magnitude Elastic Net coefficients
after nested-CV tuning ($C = 0.01$, $\alpha = 0.1$).
The strong regularisation ($C = 0.01$) shrinks coefficients toward zero, and
21 out of 26 features survive with non-zero weights.
Positive coefficients indicate features that, when larger for the Red
corner, increase the predicted probability of a Red victory.

\begin{table}[H]
\centering
\caption{Top 8 Elastic Net coefficients (sorted by $|\hat{\beta}|$)}.
\label{tab:coefs}
\begin{tabular}{lrl}
\toprule
Feature & $\hat{\beta}$ & Interpretation \\
\midrule
\texttt{AgeDiff}              & $-0.259$ & Older Red fighter $\rightarrow$ lower Red win probability \\
\texttt{AvgTDLandedDiff}      & $+0.123$ & Red takedown advantage $\rightarrow$ higher Red win prob. \\
\texttt{ReachCmsDiff}         & $+0.105$ & Red reach advantage $\rightarrow$ higher Red win prob. \\
\texttt{CurrentWinStreakDiff} & $+0.104$ & Longer Red active win streak $\rightarrow$ higher Red win prob. \\
\texttt{LossesDiff}           & $-0.087$ & More career losses for Red $\rightarrow$ lower Red win prob. \\
\texttt{WeightLbsDiff}        & $+0.080$ & Red weight advantage $\rightarrow$ higher Red win prob. \\
\texttt{LongestWinStreakDiff} & $+0.069$ & Longer Red peak win streak $\rightarrow$ higher Red win prob. \\
\texttt{DrawsDiff}            & $+0.057$ & More draws for Red $\rightarrow$ slight Red win advantage \\
\bottomrule
\end{tabular}
\end{table}

The negative age coefficient warrants attention: \texttt{AgeDiff} = Age$_\text{Red}$ $-$ Age$_\text{Blue}$, so the negative sign means older Red fighters are predicted to win less often.
Since Red fighters tend to be older (mean \texttt{AgeDiff} = 0.56 years), this
coefficient captures a pattern where relatively younger opponents outperform
statistical expectations, possibly reflecting some age related decline of
the more veteran Red fighters.

\subsection{Feature importance (Random Forest)}

Figure~\ref{fig:rf_importance} shows the top-20 feature importances from the
Random Forest, measured by mean decrease in impurity.
The top three features are \texttt{AgeDiff}, \texttt{AvgSigStrLandedDiff},
and \texttt{AvgTDLandedDiff}.
This broadly agrees with the Elastic Net coefficient magnitudes
(Table~\ref{tab:coefs}), giving confidence that these features carry genuine
signal rather than being artefacts of one particular model class.

\begin{figure}[H]
\centering
\includegraphics[width=0.78\textwidth]{../figures/feature_importance/random_forest_diff_importance.png}
\caption{Random Forest feature importances (mean decrease in impurity) for
the top 20 features. \texttt{AgeDiff} dominates, with grappling and striking
averages following.}
\label{fig:rf_importance}
\end{figure}

\subsection{Bias-variance discussion}

Logistic regression has the lowest variance but a linear decision boundary
that cannot capture threshold effects (e.g., a reach advantage may matter
very differently in the heavyweight division versus flyweight).
Random Forest and XGBoost have lower bias and can fit interactions, but their
higher variance must be controlled through regularisation
(\texttt{min\_samples\_leaf = 10} and \texttt{min\_child\_weight = 5}).
Despite the added flexibility, ensemble methods do not substantially improve
over regularised logistic regression (log loss difference $< 0.004$).
This suggests that the signal-to-noise ratio in public UFC statistics is
intrinsically low, and the variance cost of flexible models nearly cancels
their bias benefit.

The fact that $C = 0.01$ (strong regularisation) was selected by inner CV
for the Elastic Net also points in this direction: the CV is protecting
against overfitting to noise, not trying to fit a complex signal.

\subsection{Limitations}

\begin{itemize}
  \item \textbf{Bookmaker gap}: The persistent $\approx$5.5 percentage-point
  accuracy gap to the bookmaker suggests important predictors (training camp
  quality, recent injuries, travel, stylistic matchups) are unobservable
  from public statistics.
  \item \textbf{Missing data}: Many earlier bouts lack complete fighter
  statistics and median imputation introduces noise that may weaken model
  discrimination for these fights.
  \item \textbf{Label asymmetry}: The Red/Blue assignment is non-random as
  Red fighters tend to be higher-ranked and older.
  This means the difference features carry a structural bias toward the
  Red corner that the model must implicitly correct for.
  \item \textbf{Fixed hyperparameters}: Hyperparameters for Random Forest
  and XGBoost were not formally tuned.
  \item \textbf{Chronological split}: Using only the last 20\% of fights as
  the test set means test performance is measured on recent fights.
  If fighter statistics tracking or matchmaking strategy has changed over
  time, the train-test gap may reflect a distributional shift rather than
  a true generalisation error.
\end{itemize}

%% -------------------------------------------------------------------------
\section{Summary}
%% -------------------------------------------------------------------------

We framed UFC fight outcome prediction as a binary classification problem
and compared three model families on a chronological test set of 1{,}255
bouts.
All models comfortably beat the naive majority-class baseline (AUC 0.5) and
extract meaningful signal from the difference feature representation.
No model closes the gap to the bookmaker benchmark (AUC 0.743), confirming
that market-implied probabilities encode information not accessible from
public fighter statistics alone.

The Random Forest achieved the best overall performance (log loss 0.633,
AUC 0.693), but the margin over regularised logistic regression is small.
The Elastic Net with nested cross-validation provided the best
interpretability: inner cross-validation selected strong regularisation
($C = 0.01$), and 21 out of 26 features survived the Lasso step.
The most important predictors across both model families were average
takedowns landed, age difference, reach advantage, and win-streak difference.

The most practically significant finding is the consistency of the
bookmaker gap across all models.
It establishes a clear ceiling: the information contained in publicly
available cumulative career statistics is real but insufficient, and
closing the gap would likely require access to information that is not
neatly available in some public dataset.

\end{document}
